\chapter{Revisão Bibliográfica}
\label{capitulo:revisao_bibliografica}

Este capítulo apresenta a base teórica e mercadológica que fundamenta o desenvolvimento do simulador \textit{Royal Legacy}. Para compreender a arquitetura do projeto, é necessário explorar os conceitos tecnológicos aplicados na sua construção, bem como analisar softwares já existentes que resolvem problemas semelhantes.

A seção \ref{secao:fundamentacao_teorica} dedica-se à fundamentação teórica, descrevendo os conceitos de xadrez computacional, os protocolos de comunicação padronizados e as tecnologias (motores e linguagens de programação) escolhidas para o desenvolvimento. Em seguida, a seção \ref{secao:trabalhos_correlatos} apresenta uma análise de trabalhos correlatos e plataformas de xadrez amplamente utilizadas no mercado, destacando as diferenças e justificando as escolhas arquiteturais deste projeto em relação ao cenário atual.

\section{Fundamentação teórica}
\label{secao:fundamentacao_teorica}

A fundamentação teórica estabelece os conceitos e as ferramentas obtidos na literatura técnica que viabilizaram a construção do simulador. Segundo \citeonline{RaulSidneiWazlawick2021}, esta etapa é essencial para embasar cientificamente as decisões de projeto. A seguir, são detalhados os protocolos e \textit{frameworks} centrais do trabalho.

\subsection{Xadrez Computacional e Protocolos (FEN e UCI)}

O desenvolvimento de interfaces para xadrez exige a padronização da forma como o tabuleiro e as jogadas são representados pelo computador. A notação FEN (\textit{Forsyth-Edwards Notation}) é o padrão global para descrever as posições de um jogo de xadrez em uma única linha de texto. Ela mapeia quais peças estão em quais casas, de quem é o turno e as disponibilidades de roque \cite{Hooper1992}.

Além da representação do tabuleiro, a comunicação com motores de Inteligência Artificial requer uma linguagem comum. O protocolo UCI (\textit{Universal Chess Interface}) é um padrão aberto que permite que interfaces gráficas (GUI) conversem com motores de xadrez de forma independente. O motor recebe a posição atual via texto e devolve a melhor jogada calculada, garantindo a separação entre a lógica visual e a lógica matemática profunda.

\subsection{Godot Engine e GDScript}

A Godot Engine é um motor de jogo (\textit{Game Engine}) de código aberto e multiplataforma, amplamente reconhecido pela sua arquitetura flexível baseada em nós (\textit{nodes}) e cenas. Diferente de outras ferramentas do mercado, a Godot oferece um ambiente unificado para o desenvolvimento 2D e 3D, sendo altamente otimizada para a criação de interfaces de usuário (UI) complexas \cite{Linietsky2023}.

Para a manipulação da lógica interna do jogo, utilizou-se o GDScript, a linguagem de programação nativa da Godot. O GDScript possui uma sintaxe inspirada em Python, sendo dinamicamente tipada e profundamente integrada ao motor, o que facilita o acesso rápido aos elementos visuais, instanciamento de cenas e manipulação de estruturas de dados, como os Dicionários utilizados para a troca de temas no \textit{Royal Legacy}.

\subsection{Linguagem Python}

Python é uma linguagem de programação de alto nível, interpretada e de propósito geral, conhecida por sua legibilidade e vasta biblioteca padrão \cite{VanRossum2009}. Neste projeto, Python foi escolhido para atuar como um \textit{middleware} (ponte) devido à sua robustez na manipulação de processos do sistema operacional e fluxos de entrada e saída de dados. Através de bibliotecas específicas, o script Python executa o binário da Inteligência Artificial em segundo plano, recebe os comandos da interface na Godot, formata os dados e os repassa via protocolo UCI, resolvendo os problemas de travamento da interface gráfica.

\subsection{Motor de Inteligência Artificial Stockfish}

O Stockfish é atualmente um dos motores de xadrez de código aberto mais fortes do mundo. Em vez de possuir uma interface visual própria, ele é um programa executável operado inteiramente por linha de comando. Nas suas versões mais recentes, o Stockfish combina algoritmos tradicionais de busca em profundidade, como o \textit{Alpha-Beta Pruning}, com avaliações baseadas em redes neurais eficientes (NNUE), permitindo-lhe avaliar milhões de posições por segundo com precisão sobre-humana \cite{Stockfish2023}.

\section{Trabalhos correlatos}
\label{secao:trabalhos_correlatos}

Esta seção apresenta os principais softwares, plataformas e repositórios relacionados ao tema deste estudo. A análise destes sistemas permite não apenas identificar os padrões da indústria, mas também reconhecer as bases comunitárias que apoiaram o desenvolvimento arquitetural do simulador \textit{Royal Legacy}.

\subsection{Godot With Me: Chess (Repositório Open-Source)}

O projeto \textit{Godot With Me: Chess} é um repositório de código aberto hospedado no GitHub que serve como um guia prático para a implementação das regras básicas de xadrez utilizando a Godot Engine. Ele aborda a movimentação em \textit{grid} (grade) e o instanciamento das peças 2D no tabuleiro \cite{GodotWithMe2023, GodotWithMeVideo}.

Este repositório serviu como principal referência inicial para o desenvolvimento do \textit{Royal Legacy}, auxiliando no entendimento da manipulação de matrizes em GDScript e na estruturação visual das casas lógicas. No entanto, o presente trabalho diferencia-se e expande significativamente o escopo desse repositório base ao implementar três grandes melhorias estruturais: a integração de uma Inteligência Artificial externa (Stockfish) operando em segundo plano via Python; a criação de um gerenciador de cenas estruturado para menus complexos; e um sistema dinâmico de troca de temas visuais via Dicionários de dados.

\subsection{Lichess (Plataforma Web Open-Source)}

O Lichess é um dos maiores servidores de xadrez da Internet, sendo totalmente gratuito e de código aberto. A plataforma permite partidas \textit{multiplayer} e contra a máquina, utilizando o próprio motor Stockfish nos seus servidores (\textit{backend}) ou através da compilação em WebAssembly no navegador do usuário (\textit{frontend}). 

A figura \ref{figura:lichess_interface} apresenta o tabuleiro padrão da plataforma. Embora seja uma referência mundial, o Lichess possui uma arquitetura focada em tecnologias web (cliente-servidor). O presente trabalho difere por ser uma aplicação \textit{desktop} nativa focada em renderização através de uma \textit{Game Engine}, dispensando a necessidade de navegadores para a integração com a Inteligência Artificial.

%\figura{caminho}{tamanho cm}{legenda}{referencia}
\figura{./imagens/lichess.png}{12cm}{Interface de partida contra o computador na plataforma Lichess.}{figura:lichess_interface}

\subsection{Lucas Chess}

O \textit{Lucas Chess} é um software \textit{desktop} voltado para o treinamento e ensino de xadrez. O seu grande diferencial é a inclusão de dezenas de motores de IA diferentes, configurados para simular adversários de diversas faixas de força (rating ELO), de iniciantes a Grandes Mestres.

Assim como o projeto \textit{Royal Legacy}, o Lucas Chess utiliza a integração local com motores via protocolo UCI. No entanto, a sua interface gráfica é construída com ferramentas tradicionais de janelas do sistema operacional, focando mais na utilidade estatística do que na experiência audiovisual ou fluidez de animações. O uso da Godot Engine neste TCC propõe justamente uma evolução na apresentação visual e interativa desse tipo de simulador local.

\subsection{Chess.com}

O Chess.com é a plataforma comercial líder de mercado em xadrez digital. A aplicação oferece uma vasta gama de personalização de temas, avatares (bots) com personalidades distintas e análises profundas de partidas. Por ser um software proprietário, os motores lógicos utilizados por trás dos seus "bots" são fechados em servidores remotos de alto desempenho. O desenvolvimento do \textit{Royal Legacy} busca espelhar a experiência de customização visual e a seleção de dificuldades de IA encontrada no Chess.com, porém arquitetando tudo de forma local, \textit{offline} e utilizando apenas ferramentas \textit{open-source}.